\chapter[Distance-Based Methods, Decision Trees and Naive Bayes]{Distance-Based Methods,\\Decision Trees and Naive Bayes}
\label{ch:knn}
\begin{objectives}
After studying this chapter you will be able to:
\begin{itemize}
  \item compute \textbf{Euclidean, Manhattan, Minkowski, Chebyshev, Hamming} and \textbf{cosine} measures and explain why \textbf{feature scaling} matters;
  \item apply the \textbf{k-Nearest Neighbours} algorithm to classification and regression, and explain the effect of $k$;
  \item build a \textbf{decision tree} by hand using \textbf{entropy, information gain} and \textbf{Gini index};
  \item explain \textbf{overfitting in trees} and \textbf{pruning};
  \item state \textbf{Bayes' theorem} and apply the \textbf{Naive Bayes} classifier, including \textbf{Laplace smoothing} and the \textbf{Gaussian} variant;
  \item compare the three methods and choose between them;
  \item answer 2-, 5- and 10-mark university questions on these topics.
\end{itemize}
\end{objectives}

\begin{pullquote}[--- Cover \& Hart, 1967 (on nearest neighbours)]
Birds of a feather flock together: the nearest neighbour of a point carries a lot of information about its class.
\end{pullquote}

This chapter begins Unit~II (Supervised Learning). The three methods here are among the simplest and most intuitive classifiers, and each uses a different idea:

\begin{center}
\begin{tabular}{@{}lll@{}}
\toprule
\textbf{Method} & \textbf{Core idea} & \textbf{Type (Chapter~1 / 2)}\\
\midrule
k-Nearest Neighbours & ``You are like your closest neighbours'' & instance-based, non-parametric\\
Decision tree & ``Ask a sequence of yes/no questions'' & model-based, non-parametric\\
Naive Bayes & ``Use probability and Bayes' rule'' & model-based, parametric\\
\bottomrule
\end{tabular}
\end{center}

%====================================================================
\section{Distance-Based Methods}

Many learning algorithms (k-NN, K-Means, SVMs with kernels) depend on one question: \emph{how similar are two instances?} We answer it with a \textbf{distance} (small = similar) or a \textbf{similarity} (large = similar).

\begin{definition}[: Distance-based methods]
\textbf{Distance-based methods} make predictions or form groups using the \emph{distance (or similarity) between instances} in feature space: close points are assumed to be similar.
\end{definition}

\begin{definition}[: Distance metric]
A function $d(\vect{a},\vect{b})$ is a \textbf{metric} if for all points: (i)~$d(\vect{a},\vect{b})\ge0$ and $d(\vect{a},\vect{b})=0\iff\vect{a}=\vect{b}$; (ii)~symmetry $d(\vect{a},\vect{b})=d(\vect{b},\vect{a})$; (iii)~triangle inequality $d(\vect{a},\vect{c})\le d(\vect{a},\vect{b})+d(\vect{b},\vect{c})$.
\end{definition}

\subsection{Common Distance Measures}

For two $d$-dimensional points $\vect{a}=(a_1,\dots,a_d)$ and $\vect{b}=(b_1,\dots,b_d)$:

\begin{center}
\begin{tabular}{@{}p{3.3cm}p{5.8cm}p{4.4cm}@{}}
\toprule
\textbf{Measure} & \textbf{Formula} & \textbf{Use}\\
\midrule
Euclidean ($L_2$) & $\sqrt{\sum_i (a_i-b_i)^2}$ & Default for numeric data\\
Manhattan ($L_1$) & $\sum_i |a_i-b_i|$ & Grid-like paths; less sensitive to outliers\\
Minkowski ($L_p$) & $\big(\sum_i |a_i-b_i|^p\big)^{1/p}$ & General family ($p=1,2$ above)\\
Chebyshev ($L_\infty$) & $\max_i |a_i-b_i|$ & Largest single difference\\
Hamming & number of positions where $a_i\neq b_i$ & Binary / categorical strings\\
Cosine similarity & $\dfrac{\vect{a}\cdot\vect{b}}{\|\vect{a}\|\,\|\vect{b}\|}$ & Text; direction, not size\\
\bottomrule
\end{tabular}
\end{center}

\begin{example}[: Computing distances]
Let $\vect{a}=(1,2,3)$ and $\vect{b}=(4,6,3)$. The coordinate differences are $3,4,0$.
\begin{align*}
d_{L_2}&=\sqrt{3^2+4^2+0^2}=\sqrt{25}=5, &
d_{L_1}&=3+4+0=7,\\
d_{L_\infty}&=\max(3,4,0)=4, &
d_{L_3}&=(27+64+0)^{1/3}=91^{1/3}\approx4.50.
\end{align*}
Cosine similarity: $\vect{a}\cdot\vect{b}=4+12+9=25$, $\|\vect{a}\|=\sqrt{14}\approx3.742$, $\|\vect{b}\|=\sqrt{61}\approx7.810$, so $\cos\theta=25/29.22\approx0.855$.\par
Hamming: for the bit strings $10110$ and $11100$ the strings differ in positions 2 and 4, so $d_H=2$.\par
\textbf{Note the ordering} $L_\infty\le L_3\le L_2\le L_1$ for the same pair of points.
\end{example}

\begin{figure}[!htbp]
\centering
\begin{tikzpicture}[scale=1.6]
\draw[-{Stealth},gray] (-1.4,0)--(1.5,0); \draw[-{Stealth},gray] (0,-1.4)--(0,1.5);
\draw[very thick,mlred] (1,0)--(0,1)--(-1,0)--(0,-1)--cycle;
\draw[very thick,mlblue] (0,0) circle (1);
\draw[very thick,mlgreen] (-1,-1) rectangle (1,1);
\node[mlred,font=\small,fill=white,inner sep=1pt] at (0.62,0.62) {$L_1$};
\node[mlblue,font=\small,fill=white,inner sep=1pt] at (0.62,-0.82) {$L_2$};
\node[mlgreen,font=\small,fill=white,inner sep=1pt] at (1.18,0.9) {$L_\infty$};
\fill (0,0) circle (1pt) node[below left,font=\scriptsize] {origin};
\end{tikzpicture}
\caption{The set of all points at distance exactly 1 from the origin, for three metrics.}\label{fig:unitballs}
\end{figure}

\begin{explain}
Each coloured outline is the set of points that are \emph{exactly one unit away} from the origin under a different distance. Under \textbf{Euclidean} distance ($L_2$, blue) this set is a \emph{circle}. Under \textbf{Manhattan} ($L_1$, red) it is a \emph{diamond}, because you may only move along the axes. Under \textbf{Chebyshev} ($L_\infty$, green) it is a \emph{square}. The point $(1,1)$ lies on the green square but far outside the red diamond: the \emph{same pair of points} is counted as ``close'' or ``far'' depending on the metric. Hence the choice of metric changes which points are neighbours.
\end{explain}


\subsection{Feature Scaling}

Distance measures add up differences across all features, so a feature with a \emph{large numeric range} dominates the sum. We must put features on a comparable scale.

\begin{center}
\begin{tabular}{@{}p{3.6cm}p{5cm}p{4.6cm}@{}}
\toprule
\textbf{Method} & \textbf{Formula} & \textbf{Result}\\
\midrule
Min--max scaling & $x'=\dfrac{x-x_{\min}}{x_{\max}-x_{\min}}$ & Values in $[0,1]$\\[3mm]
Standardisation (z-score) & $x'=\dfrac{x-\mu}{\sigma}$ & Mean 0, standard deviation 1\\
\bottomrule
\end{tabular}
\end{center}

\begin{example}[: Why scaling matters]
Three people: $P=(\text{age }25,\ \text{salary }40000)$, $Q=(45,\ 42000)$, $R=(26,\ 80000)$. Raw Euclidean distances:
\[
d(P,Q)=\sqrt{20^2+2000^2}\approx2000.1,\qquad d(P,R)=\sqrt{1^2+40000^2}\approx40000.0 .
\]
The salary difference swamps age, so $Q$ (20 years older) looks far closer to $P$ than $R$ (1 year older) does, which is silly. Now min--max scale with age range $20$--$60$ and salary range $20\,000$--$100\,000$:
\[
P'=(0.125,\,0.250),\quad Q'=(0.625,\,0.275),\quad R'=(0.150,\,0.750).
\]
$d(P',Q')=\sqrt{0.5^2+0.025^2}\approx0.5006$ and $d(P',R')=\sqrt{0.025^2+0.5^2}\approx0.5006$: after scaling, age and salary contribute fairly.
\end{example}

\begin{warning}
Always compute the scaling parameters ($x_{\min},x_{\max},\mu,\sigma$) on the \textbf{training set only} and then apply the \emph{same} numbers to the test set. Using the test data to compute them leaks information.
\end{warning}

%====================================================================
\procon{Distance-based methods}
{\item Intuitive and easy to implement
 \item Little or no training
 \item Works for classification, regression and clustering}
{\item Very sensitive to feature scale and irrelevant features
 \item Slow for large datasets
 \item Loses meaning in high dimensions (curse of dimensionality)}
{\item k-NN and K-Means clustering
 \item Recommender systems
 \item Anomaly detection}

\section{Nearest Neighbours (k-NN)}

\subsection{The Algorithm}

k-NN is the simplest learning algorithm: it has \textbf{no training phase}. It just stores the training data (a \emph{lazy learner}) and does all the work at prediction time.

\begin{definition}[: k-Nearest Neighbours]
To predict for a new point $\vect{x}_q$:
\begin{enumerate}
  \item Compute the distance from $\vect{x}_q$ to every training point.
  \item Pick the $k$ training points with the smallest distances (the \emph{neighbourhood} $N_k(\vect{x}_q)$).
  \item \textbf{Classification:} output the \emph{majority class} among the $k$ neighbours.\\
        \textbf{Regression:} output the \emph{mean} of their target values.
\end{enumerate}
\end{definition}

\begin{example}[: k-NN classification by hand]
Training data (two features, two classes) and the query $\vect{x}_q=(5,4)$:

\begin{center}
\begin{tabular}{@{}c|cc|c|c|c@{}}
\toprule
Point & $x_1$ & $x_2$ & Class & Distance to $(5,4)$ & Rank\\
\midrule
$P_4$ & 6 & 5 & B & $\sqrt{1+1}=1.414$ & 1\\
$P_5$ & 7 & 4 & B & $\sqrt{4+0}=2.000$ & 2--3\\
$P_6$ & 5 & 6 & B & $\sqrt{0+4}=2.000$ & 2--3\\
$P_7$ & 4 & 2 & A & $\sqrt{1+4}=2.236$ & 4\\
$P_3$ & 3 & 2 & A & $\sqrt{4+4}=2.828$ & 5\\
$P_2$ & 2 & 3 & A & $\sqrt{9+1}=3.162$ & 6\\
$P_1$ & 1 & 1 & A & $\sqrt{16+9}=5.000$ & 7\\
\bottomrule
\end{tabular}
\end{center}

\begin{center}
\begin{tabular}{@{}cll@{}}
\toprule
$k$ & Neighbours & Prediction\\
\midrule
1 & $P_4$ (B) & \textbf{B}\\
3 & $P_4,P_5,P_6$ = B,B,B & \textbf{B}\\
5 & B,B,B,A,A (3 vs 2) & \textbf{B}\\
7 & all points: 4 A vs 3 B & \textbf{A}\\
\bottomrule
\end{tabular}
\end{center}
The prediction \emph{changes} from B to A when $k$ grows to 7, which shows how sensitive k-NN is to $k$.
\end{example}

\begin{figure}[!htbp]
\centering
\begin{tikzpicture}
\begin{axis}[width=8.2cm,height=8.2cm,xmin=0,xmax=8,ymin=0,ymax=8,xlabel={$x_1$},ylabel={$x_2$},
  axis equal image,grid=both,grid style={gray!20},legend style={at={(1.03,0.5)},anchor=west,font=\small}]
\addplot[only marks,mark=*,mlblue,mark size=3pt] coordinates {(1,1)(2,3)(3,2)(4,2)};
\addlegendentry{class A}
\addplot[only marks,mark=triangle*,mlorange,mark size=4pt] coordinates {(6,5)(7,4)(5,6)};
\addlegendentry{class B}
\addplot[only marks,mark=star,mlgreen,mark size=5pt,very thick] coordinates {(5,4)};
\addlegendentry{query}
\addplot[dashed,mlgreen,thick,domain=0:360,samples=100,variable=\t]({5+2.05*cos(t)},{4+2.05*sin(t)});
\addlegendentry{$k=3$}
\addplot[dotted,mlred,very thick,domain=0:360,samples=100,variable=\t]({5+2.9*cos(t)},{4+2.9*sin(t)});
\addlegendentry{$k=5$}
\end{axis}
\end{tikzpicture}
\caption{Neighbourhoods of the query point for $k=3$ and $k=5$.}\label{fig:knn}
\end{figure}

\begin{explain}
The \textbf{green star} is the query point. Blue circles are class~A and orange triangles are class~B. The \textbf{dashed green circle} is just big enough to contain the 3 nearest points: all three are orange, so $k=3$ votes B. The \textbf{dotted red circle} is large enough to contain 5 points: three orange and two blue, so $k=5$ still votes B, but by a smaller margin. If we kept enlarging the circle to include all 7 points the blue majority would win. The \emph{radius adapts} to the local density of data; this is why k-NN is called \emph{local}.
\end{explain}


\subsection{Choosing $k$}

\begin{center}
\begin{tabular}{@{}p{2.8cm}p{5.5cm}p{5.5cm}@{}}
\toprule
 & \textbf{Small $k$ (e.g.\ 1)} & \textbf{Large $k$}\\
\midrule
Decision boundary & Very jagged, follows every point & Smooth\\
Bias & Low & High\\
Variance & High (sensitive to noise) & Low\\
Risk & Overfitting & Underfitting\\
\bottomrule
\end{tabular}
\end{center}

$k$ is a \textbf{hyperparameter}, chosen by cross-validation (Chapter~2). Use an \emph{odd} $k$ for two-class problems to avoid ties. A common starting point is $k\approx\sqrt n$.

\begin{keyidea}
$1/k$ is a measure of model flexibility: $k=1$ is the \emph{most flexible} (and most variable) k-NN; $k=n$ always predicts the overall majority class (the least flexible). This is the bias--variance trade-off of Chapter~2 in action.
\end{keyidea}

\subsection{Weighted k-NN and k-NN Regression}

In \textbf{weighted k-NN}, closer neighbours get more influence, with a weight $w_i=1/d_i$ (or $1/d_i^2$).

\begin{example}[: k-NN regression]
Use the house data of Chapter~1: size $x=1,2,3,4,5$ with price $y=15,26,29,44,46$. Predict the price at $x_q=3.4$ with $k=2$. The two nearest sizes are $x=3$ (distance $0.4$, price $29$) and $x=4$ (distance $0.6$, price $44$).
\begin{itemize}
  \item \textbf{Plain average:} $\hat y=(29+44)/2=36.5$.
  \item \textbf{Weighted ($w=1/d$):} $w_1=2.5$, $w_2=1.667$; $\hat y=\dfrac{2.5(29)+1.667(44)}{2.5+1.667}=\dfrac{145.83}{4.167}=35.0$.
\end{itemize}
The weighted answer is pulled towards the closer house (price 29). For comparison, the linear model $8+8x$ from Chapter~1 predicts $35.2$.
\end{example}

\subsection{Strengths and Weaknesses}

\begin{center}
\begin{tabular}{@{}p{6.7cm}p{6.7cm}@{}}
\toprule
\textbf{Advantages} & \textbf{Disadvantages}\\
\midrule
Very simple; no training & Slow prediction: must scan all $n$ points ($O(nd)$ per query)\\
Naturally handles multiple classes & Needs a lot of memory (stores everything)\\
No assumption about the data shape & Sensitive to feature scale and irrelevant features\\
Improves as more data arrives & Suffers from the \emph{curse of dimensionality}\\
\bottomrule
\end{tabular}
\end{center}

\paragraph{Curse of dimensionality.} In high dimensions, points become almost equally far apart, so the idea of a ``nearest'' neighbour loses meaning and k-NN degrades. Remedies: feature selection and dimensionality reduction such as PCA (Chapter~8).

%====================================================================
\procon{k-Nearest neighbours (k-NN)}
{\item Simple; no training; handles many classes
 \item No assumption about the data distribution
 \item Improves as more data is added}
{\item Slow prediction and large memory
 \item Sensitive to $k$, scaling and irrelevant features
 \item Poor in high dimensions}
{\item Recommendation systems
 \item Handwriting and pattern recognition
 \item Credit rating and anomaly detection}

\section{Decision Trees}

\subsection{Idea and Structure}

A decision tree predicts by asking a sequence of simple questions about the features, like a game of ``twenty questions''.

\begin{definition}[: Decision tree]
A \textbf{decision tree} is a flowchart-like model in which each \textbf{internal node} tests one feature, each \textbf{branch} is an outcome of the test, and each \textbf{leaf} holds a prediction (a class label for \emph{classification trees}, or a number for \emph{regression trees}). The top node is the \textbf{root}.
\end{definition}

The learning problem is: \emph{which feature should we test at each node?} We choose the test that makes the resulting groups as \textbf{pure} as possible, i.e.\ each group dominated by one class.

\subsection{Measures of Impurity}

Let a node contain examples of $C$ classes with proportions $p_1,\dots,p_C$.

\begin{definition}[: Entropy]
\[
  H=-\sum_{i=1}^{C}p_i\log_2 p_i\qquad(0\log0:=0).
\]
Entropy is $0$ when the node is pure and is maximal ($\log_2C$) when all classes are equally likely. It is measured in \emph{bits}.
\end{definition}

\begin{definition}[: Gini index]
\[
  G=1-\sum_{i=1}^{C}p_i^2 .
\]
Gini is $0$ for a pure node and equals $0.5$ for a 50/50 two-class node. It is the probability that two items drawn at random from the node have different classes.
\end{definition}

\begin{figure}[!htbp]
\centering
\begin{tikzpicture}
\begin{axis}[width=10cm,height=6cm,xlabel={proportion of class 1, $p$},ylabel={impurity},
  xmin=0,xmax=1,ymin=0,ymax=1.1,domain=0.001:0.999,samples=120,axis lines=left,
  legend style={at={(0.5,1.04)},anchor=south,legend columns=3,font=\small}]
\addplot[thick,mlblue]{-x*ln(x)/ln(2)-(1-x)*ln(1-x)/ln(2)};
\addplot[thick,mlred]{2*x*(1-x)};
\addplot[thick,mlgreen,dashed]{1-max(x,1-x)};
\legend{Entropy,Gini ($\times2$ for comparison),Misclassification error}
\draw[dotted] (axis cs:0.5,0)--(axis cs:0.5,1.1);
\end{axis}
\end{tikzpicture}
\caption{Impurity measures for a two-class node as a function of $p$.}\label{fig:impurity}
\end{figure}

\begin{explain}
The horizontal axis is the \textbf{proportion of class 1} in a node: $p=0$ or $p=1$ means a \emph{pure} node and $p=0.5$ is the most mixed. All three curves are \textbf{zero at the ends} (pure) and \textbf{highest at the middle} (maximum confusion). \textbf{Entropy} (blue) is slightly more curved than \textbf{Gini} (red; drawn here as $2p(1-p)$ so that it peaks at 1 like entropy). In practice they almost always choose the same splits. The dashed green \emph{misclassification error} has a sharp corner and is less sensitive to changes, which is why it is rarely used for growing trees.
\end{explain}

\begin{definition}[: Information gain]
If a split on feature $A$ divides a set $S$ into subsets $S_v$ (one for each value $v$ of $A$), then
\[
  \text{Gain}(S,A)=H(S)-\sum_{v}\frac{|S_v|}{|S|}\,H(S_v).
\]
It is the \emph{reduction in entropy} achieved by the split. At each node the tree picks the feature with the \textbf{highest information gain}. (The ID3 algorithm does exactly this; C4.5 uses the \emph{gain ratio} to avoid favouring features with many values; CART uses the Gini index and binary splits.)
\end{definition}

\subsection{Worked Example: Building a Tree with ID3}

We want to decide whether to \textbf{play tennis} from the weather. There are 14 past days.

\begin{center}
\small
\begin{tabular}{@{}cllllc@{}}
\toprule
\textbf{Day} & \textbf{Outlook} & \textbf{Temp.} & \textbf{Humidity} & \textbf{Wind} & \textbf{Play}\\
\midrule
D1 & Sunny & Hot & High & Weak & No\\
D2 & Sunny & Hot & High & Strong & No\\
D3 & Overcast & Hot & High & Weak & Yes\\
D4 & Rain & Mild & High & Weak & Yes\\
D5 & Rain & Cool & Normal & Weak & Yes\\
D6 & Rain & Cool & Normal & Strong & No\\
D7 & Overcast & Cool & Normal & Strong & Yes\\
D8 & Sunny & Mild & High & Weak & No\\
D9 & Sunny & Cool & Normal & Weak & Yes\\
D10 & Rain & Mild & Normal & Weak & Yes\\
D11 & Sunny & Mild & Normal & Strong & Yes\\
D12 & Overcast & Mild & High & Strong & Yes\\
D13 & Overcast & Hot & Normal & Weak & Yes\\
D14 & Rain & Mild & High & Strong & No\\
\bottomrule
\end{tabular}
\end{center}

\paragraph{Step 1: entropy of the whole set.} There are 9 Yes and 5 No:
\[
H(S)=-\tfrac{9}{14}\log_2\tfrac{9}{14}-\tfrac{5}{14}\log_2\tfrac{5}{14}=0.410+0.530=\mathbf{0.940}.
\]

\paragraph{Step 2: information gain of every feature.}

\begin{center}
\begin{tabular}{@{}lccc@{}}
\toprule
\textbf{Split} & \textbf{Values (Yes, No)} & \textbf{Entropies} & \textbf{Gain}\\
\midrule
Outlook & Sunny (2,3), Overcast (4,0), Rain (3,2) & $0.971,\ 0,\ 0.971$ & $\mathbf{0.247}$\\
Humidity & High (3,4), Normal (6,1) & $0.985,\ 0.592$ & $0.152$\\
Wind & Weak (6,2), Strong (3,3) & $0.811,\ 1.000$ & $0.048$\\
Temperature & Hot (2,2), Mild (4,2), Cool (3,1) & $1.000,\ 0.918,\ 0.811$ & $0.029$\\
\bottomrule
\end{tabular}
\end{center}

\emph{Detail for Outlook:}
\[
\text{Gain}=0.940-\tfrac{5}{14}(0.971)-\tfrac{4}{14}(0)-\tfrac{5}{14}(0.971)=0.940-0.694=0.247 .
\]
\textbf{Outlook has the highest gain, so it becomes the root.}

\paragraph{Step 3: recurse on each branch.}
\begin{itemize}
  \item \textbf{Overcast:} all 4 days are Yes $\Rightarrow$ pure $\Rightarrow$ leaf \emph{Yes}.
  \item \textbf{Sunny} (D1, D2, D8, D9, D11; 2 Yes, 3 No): splitting on Humidity gives High $\to$ \{D1, D2, D8\} all No, Normal $\to$ \{D9, D11\} all Yes. Both pure, so Humidity is chosen.
  \item \textbf{Rain} (D4, D5, D6, D10, D14; 3 Yes, 2 No): splitting on Wind gives Weak $\to$ \{D4, D5, D10\} all Yes, Strong $\to$ \{D6, D14\} all No.
\end{itemize}

\begin{figure}[!htbp]
\centering
\begin{tikzpicture}[
  every node/.style={font=\small},
  dec/.style={draw,rounded corners,fill=mlblue!12,draw=mlblue,thick,minimum width=2.1cm,minimum height=0.8cm},
  yes/.style={draw,fill=mlgreen!22,draw=mlgreen,thick,minimum width=1.4cm,minimum height=0.7cm,rounded corners},
  no/.style={draw,fill=mlred!15,draw=mlred,thick,minimum width=1.4cm,minimum height=0.7cm,rounded corners},
  lab/.style={fill=white,inner sep=1.5pt,font=\footnotesize}]
\node[dec] (o) at (0,0) {Outlook?};
\node[dec] (h) at (-4.2,-2.3) {Humidity?};
\node[yes] (oc) at (0,-2.3) {Yes};
\node[dec] (w) at (4.2,-2.3) {Wind?};
\node[no]  (hn) at (-5.4,-4.6) {No};
\node[yes] (hy) at (-3.0,-4.6) {Yes};
\node[no]  (ws) at (3.0,-4.6) {No};
\node[yes] (ww) at (5.4,-4.6) {Yes};
\draw[thick] (o)--(h) node[midway,lab]{Sunny};
\draw[thick] (o)--(oc) node[midway,lab]{Overcast};
\draw[thick] (o)--(w) node[midway,lab]{Rain};
\draw[thick] (h)--(hn) node[midway,lab]{High};
\draw[thick] (h)--(hy) node[midway,lab]{Normal};
\draw[thick] (w)--(ws) node[midway,lab]{Strong};
\draw[thick] (w)--(ww) node[midway,lab]{Weak};
\end{tikzpicture}
\caption{The decision tree learned by ID3 from the Play-Tennis data.}\label{fig:tree}
\end{figure}

\begin{explain}
\textbf{Blue boxes} are decision nodes (questions), \textbf{green} and \textbf{red} boxes are leaves (final answers: Yes/No). To classify a new day, start at the \emph{root} and follow the branch matching the day's value at each question until you reach a leaf. For example, a day that is \emph{Rain} and \emph{Weak wind}: Outlook $=$ Rain $\to$ Wind? $\to$ Weak $\to$ \textbf{Yes}. The tree also reads as rules:
\textit{IF Outlook=Overcast THEN Play=Yes; IF Outlook=Sunny AND Humidity=High THEN No; \ldots}\ This is why trees are so easy to \emph{interpret}.
\end{explain}

\paragraph{The same split using the Gini index.} $G(S)=1-(\tfrac9{14})^2-(\tfrac5{14})^2=0.459$. For Outlook, $G(\text{Sunny})=G(\text{Rain})=1-(0.4)^2-(0.6)^2=0.48$ and $G(\text{Overcast})=0$, so the weighted Gini is $\tfrac5{14}(0.48)+\tfrac4{14}(0)+\tfrac5{14}(0.48)=0.343$ and the Gini gain is $0.459-0.343=0.116$, again the best feature.

\subsection{Algorithm Summary}

\begin{enumerate}
  \item Start with all training data at the root.
  \item If the node is pure (or a stopping rule applies), make it a leaf labelled with the majority class.
  \item Otherwise, evaluate every candidate split; choose the one with the highest gain (or largest drop in Gini).
  \item Split the data and \textbf{recurse} on each child.
\end{enumerate}

\paragraph{Continuous features.} Sort the values, try thresholds halfway between consecutive values (e.g.\ ``Age $\le 30$?''), and choose the threshold with the best gain.

\paragraph{Regression trees.} Leaves store the \emph{mean} target of their training points; the best split minimises the sum of squared errors (variance) in the two children.

\begin{figure}[!htbp]
\centering
\begin{tikzpicture}
\begin{axis}[width=8cm,height=6.4cm,xmin=0,xmax=10,ymin=0,ymax=10,xlabel={$x_1$},ylabel={$x_2$},ticks=none,
  legend style={at={(1.03,0.5)},anchor=west,font=\small}]
\fill[mlblue!8] (axis cs:0,0) rectangle (axis cs:4,10);
\fill[mlorange!10] (axis cs:4,0) rectangle (axis cs:10,5);
\fill[mlblue!8] (axis cs:4,5) rectangle (axis cs:10,10);
\addplot[only marks,mark=*,mlblue,mark size=2.5pt] coordinates {(1,2)(2,6)(3,4)(1.5,8)(6,8)(8,7)(7,9)};
\addlegendentry{class A}
\addplot[only marks,mark=triangle*,mlorange,mark size=3.5pt] coordinates {(5,2)(6,3)(8,1.5)(7,4)(9,3)};
\addlegendentry{class B}
\draw[very thick,mlgreen] (axis cs:4,0)--(axis cs:4,10);
\draw[very thick,mlred] (axis cs:4,5)--(axis cs:10,5);
\node[font=\scriptsize,mlgreen!50!black,fill=white,inner sep=1pt] at (axis cs:4,9.5) {split 1: $x_1<4$};
\node[font=\scriptsize,mlred,fill=white,inner sep=1pt] at (axis cs:8.2,5.4) {split 2: $x_2<5$};
\end{axis}
\end{tikzpicture}
\caption{A decision tree with two splits carves the feature space into axis-parallel rectangles.}\label{fig:partition}
\end{figure}

\begin{explain}
Every split in a tree tests \emph{one feature against a threshold}, which draws a \textbf{straight line parallel to an axis}. The \textbf{green vertical line} is the root test $x_1<4$: everything to its left is class~A. On the right, the \textbf{red horizontal line} tests $x_2<5$: below it is class~B (orange), above it class~A. The final result is three \emph{rectangular regions}, each assigned one class. Trees therefore produce \emph{boxy} decision boundaries, which suit data with axis-aligned structure but approximate diagonal boundaries only with many small steps.
\end{explain}


\subsection{Overfitting and Pruning}

A tree allowed to grow until every leaf is pure will memorise the training data, noise included (low bias, \textbf{high variance}; see Chapter~2). We control complexity by \textbf{pruning}.

\begin{center}
\begin{tabular}{@{}p{3.3cm}p{10.2cm}@{}}
\toprule
\textbf{Pre-pruning} (early stopping) & Stop growing early: limit the depth (\texttt{max\_depth}), require a minimum number of samples to split (\texttt{min\_samples\_split}) or in a leaf (\texttt{min\_samples\_leaf}), or require a minimum gain.\\
\textbf{Post-pruning} & Grow the full tree, then \emph{cut back} branches that do not improve accuracy on a validation set (reduced-error pruning) or minimise cost $=\text{error}+\alpha\cdot(\text{number of leaves})$ (cost-complexity pruning).\\
\bottomrule
\end{tabular}
\end{center}

\subsection{Strengths and Weaknesses}

\begin{center}
\begin{tabular}{@{}p{6.7cm}p{6.7cm}@{}}
\toprule
\textbf{Advantages} & \textbf{Disadvantages}\\
\midrule
Easy to understand and visualise; human-readable rules & Prone to overfitting\\
Handles numeric and categorical features & Unstable: small data changes can give a very different tree\\
No feature scaling needed & Axis-parallel (boxy) boundaries only\\
Fast prediction; performs feature selection implicitly & Biased towards features with many values\\
\bottomrule
\end{tabular}
\end{center}

The instability of single trees is fixed by combining many of them into \emph{ensembles} such as Random Forests (Chapter~5).

%====================================================================
\procon{Decision trees}
{\item Easy to understand, visualise and convert to rules
 \item Handles numeric and categorical features; no scaling needed
 \item Fast prediction}
{\item Prone to overfitting; needs pruning
 \item Unstable: small data changes can change the tree
 \item Only axis-parallel (boxy) boundaries}
{\item Credit approval and medical diagnosis
 \item Customer churn prediction
 \item Rule extraction; base learner in Random Forests}

\section{Naive Bayes}

\subsection{Bayes' Theorem}

\begin{definition}[: Bayes' theorem]
For a class $C$ and observed evidence $\vect{x}$,
\[
  P(C\mid\vect{x})=\frac{P(\vect{x}\mid C)\,P(C)}{P(\vect{x})},\qquad
  P(\vect{x})=\sum_{c}P(\vect{x}\mid c)P(c).
\]
Here $P(C\mid\vect{x})$ is the \textbf{posterior}, $P(\vect{x}\mid C)$ the \textbf{likelihood}, $P(C)$ the \textbf{prior} and $P(\vect{x})$ the \textbf{evidence}.
\end{definition}

It follows directly from the definition of conditional probability, $P(C\mid\vect{x})P(\vect{x})=P(\vect{x},C)=P(\vect{x}\mid C)P(C)$. In words: \emph{posterior} (what we believe after seeing the evidence) $\propto$ \emph{likelihood} (how well the class explains the evidence) $\times$ \emph{prior} (what we believed before).

\begin{example}[: Spam and the word ``offer'']
Suppose 30\% of all e-mail is spam, so $P(\text{spam})=0.3$ and $P(\text{ham})=0.7$. The word ``offer'' occurs in 60\% of spam, $P(\text{offer}\mid\text{spam})=0.6$, but only in 10\% of ham, $P(\text{offer}\mid\text{ham})=0.1$. An e-mail contains ``offer''. How likely is it to be spam?
\[
P(\text{spam}\mid\text{offer})=\frac{0.6\times0.3}{0.6\times0.3+0.1\times0.7}=\frac{0.18}{0.18+0.07}=\frac{0.18}{0.25}=\mathbf{0.72}.
\]
\end{example}

\begin{figure}[!htbp]
\centering
\begin{tikzpicture}[
  every node/.style={font=\small},
  lvl/.style={draw,rounded corners,thick,minimum height=0.8cm,inner sep=4pt},
  lab/.style={fill=white,inner sep=1.5pt,font=\footnotesize}]
\node[lvl,fill=gray!12] (r) at (0,0) {All e-mails};
\node[lvl,fill=mlred!15,draw=mlred] (s) at (4.2,1.6) {Spam};
\node[lvl,fill=mlgreen!15,draw=mlgreen] (h) at (4.2,-1.6) {Ham};
\node[lvl,fill=mlred!25,draw=mlred] (so) at (8.6,2.3) {contains ``offer''};
\node[lvl,fill=gray!8] (sn) at (8.6,0.9) {does not};
\node[lvl,fill=mlgreen!25,draw=mlgreen] (ho) at (8.6,-0.9) {contains ``offer''};
\node[lvl,fill=gray!8] (hn) at (8.6,-2.3) {does not};
\draw[thick] (r)--(s) node[midway,lab]{$0.3$};
\draw[thick] (r)--(h) node[midway,lab]{$0.7$};
\draw[thick] (s)--(so) node[midway,lab]{$0.6$};
\draw[thick] (s)--(sn) node[midway,lab]{$0.4$};
\draw[thick] (h)--(ho) node[midway,lab]{$0.1$};
\draw[thick] (h)--(hn) node[midway,lab]{$0.9$};
\node[font=\footnotesize,anchor=west,mlred] at (10.9,2.3) {$0.3\times0.6=0.18$};
\node[font=\footnotesize,anchor=west,mlgreen!50!black] at (10.9,-0.9) {$0.7\times0.1=0.07$};
\end{tikzpicture}
\caption{Probability tree for the spam example. Only the two ``offer'' leaves matter once we know the word appeared.}\label{fig:bayestree}
\end{figure}

\begin{explain}
Read the tree \textbf{left to right}, \emph{multiplying probabilities along a path}. First the e-mail is spam (0.3) or ham (0.7). Then, given its class, it contains ``offer'' with probability 0.6 (spam) or 0.1 (ham). The \textbf{two paths that end in ``contains offer''} have joint probabilities $0.18$ (spam) and $0.07$ (ham). Because we \emph{know} the word appeared, we only look at those two paths, and the posterior is the spam share of their total: $0.18/(0.18+0.07)=0.72$. The prior of 30\% has been \textbf{updated} to 72\% by the evidence.
\end{explain}

\subsection{The Naive Bayes Classifier}


\begin{definition}[: Naive Bayes classifier]
\textbf{Naive Bayes} is a \emph{supervised, probabilistic} classifier based on \textbf{Bayes' theorem} with the ``naive'' assumption that features are \emph{conditionally independent} given the class: $\hat y=\argmax_c P(c)\prod_jP(x_j\mid c)$.
\end{definition}


With many features $\vect{x}=(x_1,\dots,x_d)$, the likelihood $P(x_1,\dots,x_d\mid C)$ would need an enormous number of probabilities. Naive Bayes makes one simplifying assumption.

\begin{definition}[: Naive (conditional independence) assumption]
Given the class, the features are \textbf{independent}:
\[
  P(x_1,\dots,x_d\mid C)=\prod_{j=1}^{d}P(x_j\mid C).
\]
\end{definition}

The classification rule becomes (the denominator $P(\vect{x})$ is the same for every class, so it can be dropped):
\begin{equation}\label{eq:nb}
  \hat y=\argmax_{c}\; P(c)\prod_{j=1}^{d}P(x_j\mid c).
\end{equation}
The assumption is usually \emph{false} (in e-mail, ``free'' and ``offer'' tend to appear together), hence ``naive'', yet the classifier works surprisingly well, because we only need the class with the \emph{largest} score, not exact probabilities.

\paragraph{Training = counting.} Estimate $P(c)=N_c/N$ and $P(x_j=v\mid c)=\dfrac{\text{count}(x_j=v,\ \text{class }c)}{N_c}$ from the training data.

\subsection{Worked Example: Play-Tennis Again}

Classify the new day $\vect{x}=(\text{Sunny},\ \text{Cool},\ \text{High},\ \text{Strong})$ using the 14-day table.

\textbf{Priors:} $P(\text{Yes})=\tfrac9{14}$, $P(\text{No})=\tfrac5{14}$.

\begin{center}
\begin{tabular}{@{}lcc@{}}
\toprule
\textbf{Likelihood} & \textbf{Given Yes (9 days)} & \textbf{Given No (5 days)}\\
\midrule
$P(\text{Sunny}\mid\cdot)$ & $2/9$ & $3/5$\\
$P(\text{Cool}\mid\cdot)$ & $3/9$ & $1/5$\\
$P(\text{High}\mid\cdot)$ & $3/9$ & $4/5$\\
$P(\text{Strong}\mid\cdot)$ & $3/9$ & $3/5$\\
\bottomrule
\end{tabular}
\end{center}

\[
\text{score(Yes)}=\tfrac{9}{14}\cdot\tfrac29\cdot\tfrac39\cdot\tfrac39\cdot\tfrac39=0.00529,\qquad
\text{score(No)}=\tfrac{5}{14}\cdot\tfrac35\cdot\tfrac15\cdot\tfrac45\cdot\tfrac35=0.02057 .
\]
\textbf{Decision:} $0.02057>0.00529\Rightarrow$ predict \textbf{No}.\\
\textbf{Posterior probabilities} (normalise): $P(\text{No}\mid\vect{x})=\dfrac{0.02057}{0.02057+0.00529}=0.795$ and $P(\text{Yes}\mid\vect{x})=0.205$.


\subsection{Laplace Smoothing}

\begin{warning}
\textbf{The zero-frequency problem.} If a value never occurs with a class in training (e.g.\ ``Overcast'' never occurs with ``No''), then $P(\text{Overcast}\mid\text{No})=0$ and the \emph{whole product becomes zero}, whatever the other features say.
\end{warning}

\textbf{Laplace (add-one) smoothing} adds a pseudo-count $\alpha$ (usually $1$) to every count:
\[
  P(x_j=v\mid c)=\frac{\text{count}(x_j=v,c)+\alpha}{N_c+\alpha V_j},
\]
where $V_j$ is the number of different values of feature $j$.

\begin{example}[: Smoothing]
$\text{count}(\text{Overcast},\text{No})=0$, $N_{\text{No}}=5$, and Outlook has $V=3$ values. With $\alpha=1$: $P(\text{Overcast}\mid\text{No})=\dfrac{0+1}{5+3}=\dfrac18=0.125$ instead of $0$.
\end{example}

\subsection{Variants of Naive Bayes}

\begin{center}
\begin{tabular}{@{}p{3.4cm}p{4.6cm}p{5.2cm}@{}}
\toprule
\textbf{Variant} & \textbf{Feature type} & \textbf{Model for $P(x_j\mid c)$}\\
\midrule
Gaussian & Continuous & Normal density with class-wise mean $\mu_{jc}$ and variance $\sigma_{jc}^2$\\
Multinomial & Counts (word counts) & Relative frequency of word in class (with smoothing)\\
Bernoulli & Binary (word present/absent) & Probability of presence/absence\\
\bottomrule
\end{tabular}
\end{center}

For the \textbf{Gaussian} variant,
\[
  P(x_j\mid c)=\frac{1}{\sqrt{2\pi}\,\sigma_{jc}}\exp\!\left(-\frac{(x_j-\mu_{jc})^2}{2\sigma_{jc}^2}\right).
\]

\begin{example}[: Gaussian Naive Bayes with one feature]
Heights: males have $\mu=175$, $\sigma=7$ cm; females have $\mu=162$, $\sigma=6$ cm; priors are equal. Classify a person of height $170$ cm.
\begin{align*}
P(170\mid\text{male})&=\tfrac{1}{7\sqrt{2\pi}}e^{-25/98}=0.0570\times0.775=0.0442,\\
P(170\mid\text{female})&=\tfrac{1}{6\sqrt{2\pi}}e^{-64/72}=0.0665\times0.411=0.0273.
\end{align*}
Since the priors are equal, the larger likelihood wins: \textbf{male}, with posterior $0.0442/(0.0442+0.0273)\approx0.62$. (Only 62\% confident because $170$ is plausible for both.)
\end{example}

\begin{keyidea}
\textbf{Numerical tip.} Multiplying many small probabilities underflows to $0$. Real implementations add \textbf{logarithms} instead: $\hat y=\argmax_c\big[\ln P(c)+\sum_j\ln P(x_j\mid c)\big]$.
\end{keyidea}

\begin{center}
\begin{tabular}{@{}p{6.7cm}p{6.7cm}@{}}
\toprule
\textbf{Advantages} & \textbf{Disadvantages}\\
\midrule
Very fast to train and predict (just counting) & Independence assumption rarely true\\
Works well with little data and many features & Probability estimates can be poorly calibrated\\
Excellent baseline for text (spam, sentiment) & Zero-frequency problem (needs smoothing)\\
Handles multiple classes naturally & Cannot learn interactions between features\\
\bottomrule
\end{tabular}
\end{center}

%====================================================================
\procon{Naive Bayes}
{\item Very fast to train and predict
 \item Works well with little data and many features
 \item Excellent baseline for text classification}
{\item Independence assumption is rarely true
 \item Probability estimates can be poorly calibrated
 \item Zero-frequency problem (needs smoothing)}
{\item Spam filtering
 \item Sentiment analysis
 \item Document and news classification
 \item Medical diagnosis}

\section{Putting It Together: The Three Methods in Python}

\begin{tryit}
Compare all three on the classic \emph{Iris} flower dataset (150 flowers, 4 measurements, 3 species).
\begin{lstlisting}
from sklearn.datasets import load_iris
from sklearn.model_selection import train_test_split, cross_val_score
from sklearn.preprocessing import StandardScaler
from sklearn.pipeline import make_pipeline
from sklearn.neighbors import KNeighborsClassifier
from sklearn.tree import DecisionTreeClassifier
from sklearn.naive_bayes import GaussianNB

X, y = load_iris(return_X_y=True)
Xtr, Xte, ytr, yte = train_test_split(
    X, y, test_size=0.3, random_state=42, stratify=y)

models = {
  "k-NN (k=5)": make_pipeline(StandardScaler(),
                              KNeighborsClassifier(n_neighbors=5)),
  "Decision tree": DecisionTreeClassifier(max_depth=3, random_state=0),
  "Naive Bayes": GaussianNB(),
}
for name, m in models.items():
    m.fit(Xtr, ytr)
    cv = cross_val_score(m, X, y, cv=5).mean()
    print(f"{name:14s} test = {m.score(Xte, yte):.3f}  CV = {cv:.3f}")
\end{lstlisting}
\textbf{Output:}
\begin{center}
\begin{tabular}{@{}lcc@{}}
\toprule
\textbf{Model} & \textbf{Test accuracy} & \textbf{5-fold CV accuracy}\\
\midrule
k-NN ($k=5$) & 0.911 & 0.960\\
Decision tree (depth 3) & 0.933 & 0.960\\
Naive Bayes & 0.911 & 0.953\\
\bottomrule
\end{tabular}
\end{center}
All three do well here. The single test split (45 flowers) is noisy, so the \emph{cross-validated} numbers are more reliable and show the methods are practically tied.

\medskip\noindent\textbf{Effect of $k$} (same data, scaled):
\begin{center}
\begin{tabular}{@{}cccl@{}}
\toprule
$k$ & Train acc. & Test acc. & Comment\\
\midrule
1 & 1.000 & 0.933 & memorises training data\\
3 & 0.971 & 0.911 & \\
5 & 0.981 & 0.911 & \\
15 & 0.981 & 0.933 & smooth boundary\\
51 & 0.876 & 0.822 & \textbf{underfits}: both fall\\
\bottomrule
\end{tabular}
\end{center}
A huge $k$ averages over too much of the data and accuracy drops on \emph{both} sets: underfitting, exactly as predicted.

\medskip\noindent\textbf{A tiny tree you can read} (\texttt{max\_depth=2}, printed with \texttt{export\_text}):
\begin{lstlisting}[numbers=none]
|--- petal_wid <= 0.70            -> class 0 (setosa)
|--- petal_wid >  0.70
|   |--- petal_wid <= 1.55       -> class 1 (versicolor)
|   |--- petal_wid >  1.55       -> class 2 (virginica)
\end{lstlisting}
Just \emph{two questions about petal width} separate the three species.
\end{tryit}

\section{Comparison of the Three Methods}

\begin{center}
\small
\begin{tabular}{@{}p{3.1cm}p{3.5cm}p{3.5cm}p{3.5cm}@{}}
\toprule
 & \textbf{k-NN} & \textbf{Decision tree} & \textbf{Naive Bayes}\\
\midrule
Learner type & Lazy, instance-based & Eager, model-based & Eager, model-based\\
Training cost & None & Moderate & Very low\\
Prediction cost & High & Very low & Very low\\
Assumption & Similar points share classes & Axis-aligned splits & Features independent given class\\
Needs scaling? & \textbf{Yes} & No & No (Gaussian: no)\\
Interpretable? & Low & \textbf{High} & Moderate\\
Main hyperparameter & $k$, metric & depth, min samples & smoothing $\alpha$\\
Typical weakness & Slow; high dimensions & Overfits; unstable & Independence assumption\\
\bottomrule
\end{tabular}
\end{center}

%====================================================================
\begin{summary}
\begin{itemize}
  \item Distance measures ($L_1,L_2,L_p,L_\infty$, Hamming, cosine) define similarity; \textbf{scale features} first (min--max or z-score, fitted on training data only).
  \item \textbf{k-NN} stores data and predicts by majority vote (classification) or mean (regression) of the $k$ closest points. Small $k$ $\Rightarrow$ high variance; large $k$ $\Rightarrow$ high bias.
  \item \textbf{Decision trees} split on the feature with highest \textbf{information gain} (entropy) or largest Gini reduction. Full trees overfit; control with \textbf{pruning}.
  \item Entropy $H=-\sum p_i\log_2p_i$; Gini $=1-\sum p_i^2$; Gain $=H(S)-\sum\frac{|S_v|}{|S|}H(S_v)$.
  \item \textbf{Naive Bayes}: $\hat y=\argmax_c P(c)\prod_jP(x_j\mid c)$, assuming conditional independence; use \textbf{Laplace smoothing} and log-probabilities; variants: Gaussian, multinomial, Bernoulli.
  \item In the Play-Tennis example the tree root is Outlook (gain $0.247$) and Naive Bayes predicts ``No'' with posterior $0.795$ for (Sunny, Cool, High, Strong).
\end{itemize}
\end{summary}

%====================================================================
\section*{Exam Corner}
\addcontentsline{toc}{section}{Exam Corner}

\subsection*{A.\ Very Short Answer Questions (2 marks each)}
\begin{enumerate}
  \item Define Euclidean and Manhattan distance. \qmarks{2}
  \item Why is feature scaling necessary for k-NN? \qmarks{2}
  \item Why is k-NN called a lazy learner? \qmarks{2}
  \item What happens when $k=1$ and when $k=n$ in k-NN? \qmarks{2}
  \item Define entropy and Gini index. \qmarks{2}
  \item What is information gain? \qmarks{2}
  \item What is pruning? Name its two types. \qmarks{2}
  \item State Bayes' theorem. \qmarks{2}
  \item Why is Naive Bayes called ``naive''? \qmarks{2}
  \item What is Laplace smoothing and why is it needed? \qmarks{2}
\end{enumerate}

\begin{modelans}[: Q4 and Q10]
\textbf{Q4.} For $k=1$ the prediction is the label of the single nearest point: the boundary is very jagged, bias is low, variance is high (overfitting). For $k=n$ every query gets the overall majority class: the model is a constant, bias is high and variance is zero (underfitting).\par
\textbf{Q10.} Laplace smoothing adds a pseudo-count (usually 1) to every feature-value count so that $P(x_j\mid c)=\frac{\text{count}+1}{N_c+V_j}$. It is needed to avoid zero probabilities for unseen feature--class combinations, which would force the whole product $\prod P(x_j\mid c)$ to zero.
\end{modelans}

\subsection*{B.\ Short Answer Questions (5 marks each)}
\begin{enumerate}
  \item Explain the k-NN algorithm for classification with an example. \qmarks{5}
  \item Discuss the effect of the value of $k$ on k-NN (bias, variance, over/underfitting). \qmarks{5}
  \item List the advantages and disadvantages of k-NN. What is the curse of dimensionality? \qmarks{5}
  \item Explain the structure of a decision tree. Differentiate entropy and Gini index. \qmarks{5}
  \item Explain overfitting in decision trees. How can it be reduced? \qmarks{5}
  \item Explain the Naive Bayes classifier and its independence assumption. \qmarks{5}
  \item Compare k-NN, decision trees and Naive Bayes in a table. \qmarks{5}
\end{enumerate}

\begin{modelans}[: B4 (outline for 5 marks)]
\emph{Structure:} root node (first test), internal nodes (feature tests), branches (outcomes), leaves (class label or value); classification of a new point follows the path from root to a leaf.\par
\emph{Entropy} $H=-\sum p_i\log_2p_i$: measures disorder of class labels; $0$ for a pure node, $1$ for a 50/50 binary node.
\emph{Gini} $=1-\sum p_i^2$: probability of mis-classifying a random item; $0$ for pure, $0.5$ for 50/50.\par
\begin{center}\small
\begin{tabular}{@{}lll@{}}
\toprule
 & \textbf{Entropy} & \textbf{Gini}\\
\midrule
Formula & $-\sum p_i\log_2 p_i$ & $1-\sum p_i^2$\\
Range (2 classes) & $0$ to $1$ & $0$ to $0.5$\\
Cost & Needs logarithm & Cheaper (no log)\\
Used in & ID3, C4.5 & CART\\
\bottomrule
\end{tabular}
\end{center}
\end{modelans}

\subsection*{C.\ Long Answer Questions (10 marks each)}
\begin{enumerate}
  \item Explain the ID3 algorithm. Using the 14-day Play-Tennis table, compute the information gain of \emph{Outlook} and show that it is chosen as the root. \qmarks{10}
  \item Explain the Naive Bayes classifier with Laplace smoothing. Use the Play-Tennis data to classify (Sunny, Cool, High, Strong). \qmarks{10}
  \item Describe the k-NN algorithm for classification and regression. Illustrate with a numerical example and discuss the choice of $k$ and distance metric. \qmarks{10}
  \item Explain distance measures used in machine learning with examples. Why must features be scaled? \qmarks{10}
  \item Discuss pre-pruning and post-pruning in decision trees. How do trees compare to k-NN and Naive Bayes in terms of interpretability and cost? \qmarks{10}
\end{enumerate}

\begin{modelans}[: C1 (key steps)]
1.~$H(S)=0.940$ (9 Yes, 5 No).\quad
2.~Outlook: Sunny (2Y,3N) $H=0.971$; Overcast (4Y,0N) $H=0$; Rain (3Y,2N) $H=0.971$.\quad
3.~$\text{Gain}=0.940-\tfrac5{14}(0.971)-\tfrac4{14}(0)-\tfrac5{14}(0.971)=0.247$.\quad
4.~Gains of the others: Humidity $0.152$, Wind $0.048$, Temperature $0.029$. Outlook is largest, so it is the root.\quad
5.~Overcast is pure (Yes); Sunny splits on Humidity; Rain splits on Wind. Draw the tree of Figure~\ref{fig:tree}.
\end{modelans}

\subsection*{D.\ Numerical Problems}
\begin{enumerate}
  \item Find the Euclidean and Manhattan distance between $(2,3)$ and $(6,6)$. \textit{(Answer: $5$ and $7$.)}
  \item Training points: $(1,2)$ A, $(2,1)$ A, $(3,3)$ A, $(4,5)$ B, $(5,4)$ B. Classify the query $(4,4)$ using $k=3$ and $k=5$ (Euclidean). \textit{(Distances: $1,1,1.414,3.606,3.606$. Answer: $k=3$: B,B,A $\to$ \textbf{B}; $k=5$: 2B vs 3A $\to$ \textbf{A}.)}
  \item Distances to a query are $1$ and $2$ with target values $10$ and $20$. Give the weighted ($1/d$) k-NN regression value for $k=2$. \textit{(Answer: $\frac{1\cdot10+0.5\cdot20}{1.5}=13.33$.)}
  \item Compute the entropy of a node with 6 positive and 2 negative examples. \textit{(Answer: $-0.75\log_2 0.75-0.25\log_2 0.25=0.811$.)}
  \item A node with 4 positives and 4 negatives is split into $(3+,1-)$ and $(1+,3-)$. Compute the information gain. \textit{(Answer: $1-\tfrac48(0.811)-\tfrac48(0.811)=0.189$.)}
  \item Compute the Gini index of a node with 3 ``yes'' and 1 ``no''. \textit{(Answer: $1-0.5625-0.0625=0.375$.)}
  \item $P(\text{spam})=0.4$, $P(\text{``free''}\mid\text{spam})=0.5$, $P(\text{``free''}\mid\text{ham})=0.05$. Find $P(\text{spam}\mid\text{``free''})$. \textit{(Answer: $\frac{0.4\cdot0.5}{0.4\cdot0.5+0.6\cdot0.05}=\frac{0.20}{0.23}=0.870$.)}
  \item A word never appears in the 100 words of class ``ham''; vocabulary size is 50. Find its smoothed probability. \textit{(Answer: $\frac{0+1}{100+50}=0.0067$.)}
\end{enumerate}

\subsection*{E.\ Multiple-Choice Questions}
\begin{enumerate}
  \item k-NN does its main computation during: \\
  (a) training \quad (b) prediction \quad (c) pruning \quad (d) scaling
  \item Which distance is the sum of absolute coordinate differences? \\
  (a) Euclidean \quad (b) Chebyshev \quad (c) Manhattan \quad (d) Cosine
  \item The entropy of a pure node is: \\
  (a) 1 \quad (b) 0.5 \quad (c) 0 \quad (d) $\infty$
  \item A decision tree chooses the split with the: \\
  (a) lowest gain \quad (b) highest information gain \quad (c) most branches \quad (d) fewest samples
  \item Which method needs feature scaling the most? \\
  (a) Decision tree \quad (b) Naive Bayes \quad (c) k-NN \quad (d) none
  \item Overfitting of a decision tree is reduced by: \\
  (a) growing it deeper \quad (b) pruning \quad (c) removing the root \quad (d) adding duplicate data
  \item In Naive Bayes the features are assumed to be: \\
  (a) identical \quad (b) correlated \quad (c) conditionally independent given the class \quad (d) Gaussian only
  \item Laplace smoothing prevents: \\
  (a) overfitting of trees \quad (b) zero probabilities \quad (c) scaling errors \quad (d) ties in k-NN
  \item The Gini index of a 50/50 two-class node is: \\
  (a) 1 \quad (b) 0.5 \quad (c) 0.25 \quad (d) 0
  \item A very large $k$ in k-NN typically causes: \\
  (a) overfitting \quad (b) underfitting \quad (c) zero error \quad (d) no change
\end{enumerate}
\textbf{Answers:} 1-b,\ 2-c,\ 3-c,\ 4-b,\ 5-c,\ 6-b,\ 7-c,\ 8-b,\ 9-b,\ 10-b.

\subsection*{F.\ Practice Problems}
\begin{enumerate}
  \item Using the Play-Tennis table, find the best feature to split on inside the \emph{Sunny} branch using information gain, and verify that Humidity gives gain $0.971$.
  \item Classify (Overcast, Mild, High, Weak) with Naive Bayes using add-one smoothing.
  \item \textbf{Coding:} using \texttt{sklearn}, plot training and test accuracy of k-NN on Iris for $k=1,3,\dots,31$ and identify the best $k$ by 5-fold cross-validation.
  \item \textbf{Coding:} fit \texttt{DecisionTreeClassifier} with \texttt{max\_depth} from 1 to 10 and plot train vs.\ test accuracy. At which depth does overfitting begin?
\end{enumerate}

\begin{revision}
\begin{center}\small
\begin{tabular}{@{}p{3.6cm}p{9.6cm}@{}}
\toprule
\textbf{Topic} & \textbf{One-line recall}\\
\midrule
Euclidean / Manhattan & $\sqrt{\sum(a_i-b_i)^2}$ / $\sum|a_i-b_i|$.\\
Scaling & Min--max $\frac{x-x_{\min}}{x_{\max}-x_{\min}}$; z-score $\frac{x-\mu}{\sigma}$; fit on train only.\\
k-NN & Lazy; majority vote / mean of $k$ closest; small $k$ overfits, large $k$ underfits.\\
Weighted k-NN & $w_i=1/d_i$; closer neighbours count more.\\
Entropy / Gini & $-\sum p\log_2p$ / $1-\sum p^2$; both $0$ if pure.\\
Information gain & $H(S)-\sum\frac{|S_v|}{|S|}H(S_v)$; choose the maximum.\\
Pruning & Pre (max depth, min samples) / post (cut back, cost-complexity).\\
Tree boundary & Axis-parallel rectangles.\\
Bayes' theorem & Posterior $=$ likelihood $\times$ prior $/$ evidence.\\
Naive Bayes & $\argmax_cP(c)\prod_jP(x_j\mid c)$; conditional independence; Laplace $\frac{n+1}{N_c+V}$.\\
Play-Tennis results & Root Outlook (gain $0.247$); NB for (S,C,H,Str): No, $P=0.795$.\\
\bottomrule
\end{tabular}
\end{center}
\end{revision}
